\documentclass[a4paper,9pt]{article}
\usepackage[utf8]{inputenc}
\usepackage[T1]{fontenc}
\usepackage[ngerman]{babel}
\usepackage{listings}
\usepackage{booktabs}
\usepackage{xcolor}
\usepackage{tcolorbox}
\usepackage{enumitem}
\usepackage[margin=1.2cm]{geometry}

\setlist{noitemsep, topsep=1pt, parsep=0pt}
\setlength{\parskip}{1pt}
\setlength{\parindent}{0pt}

\tcbset{
	colback=gray!10, colframe=gray!40,
	boxsep=1pt, top=1pt, bottom=1pt, left=2pt, right=2pt,
	fontupper=\small
}

\definecolor{codebg}{rgb}{0.95,0.95,0.95}
\definecolor{commentcolor}{rgb}{0.4,0.4,0.4}
\definecolor{keywordcolor}{rgb}{0.0,0.3,0.7}
\definecolor{stringcolor}{rgb}{0.6,0.1,0.1}

\lstset{
	language=C,
	backgroundcolor=\color{codebg},
	basicstyle=\ttfamily\footnotesize,
	keywordstyle=\color{keywordcolor}\bfseries,
	commentstyle=\color{commentcolor}\itshape,
	stringstyle=\color{stringcolor},
	showstringspaces=false,
	frame=single, framerule=0pt,
	rulecolor=\color{codebg},
	breaklines=true,
	tabsize=2,
	aboveskip=2pt, belowskip=2pt,
	escapeinside={(*@}{@*)},
}

\begin{document}
	\tableofcontents
	\newpage
	
	%---------------------------------------------------------------
	\section{Dateien}
	Nur Byte-Streams; weitere Organisation Aufgabe des Programms.
	
	\textbf{Textdatei:} Zeichencodes (UTF-8, Latin-1, ASCII); C unterstützt eigentlich nur ASCII.\\
	\textbf{Binärdatei:} Raw bytes, keine Zeichenkodierung; Programm muss Aufbau kennen (Magic Numbers).
	
	\subsubsection{fopen}
	\begin{lstlisting}
		FILE *fd = fopen(path, "r");
		if (!fd) perror("fopen");
		// Modi: r, w, a, rb, wb
	\end{lstlisting}
	
	\subsubsection{fprintf / fscanf}
	\begin{lstlisting}
		// int fprintf(FILE *stream, const char *fmt, ...);
		fprintf(fp, "Hello %s\n", name);  // ret: Anzahl Zeichen
		
		// int fscanf(FILE *stream, const char *fmt, ...);
		fscanf(f, "%d", &x);  // besser: fgets + strtok
	\end{lstlisting}
	
	\subsubsection{fputc / fgetc / fputs / fgets}
	\begin{lstlisting}
		int fputc(int c, FILE *stream); // einzelnes Zeichen schreiben
		int fgetc(FILE *stream); // einzelnes Zeichen lesen
		int fputs(const char *s, FILE *stream); // ganzen String schreiben
		char *fgets(char *s, int n, FILE *stream); // Zeichen bis '\n' in String-Buffer lesen, maximal num-1 Zeichen
		
		// typisches Lesen mit strtok:
		while (fgets(buf, STR_LEN, fd) != NULL) {
			char *tok = strtok(buf, ";\n");
			strcpy(list[i].name, tok);
			tok = strtok(NULL, ";\n");
			list[i].rang = atoi(tok);
		}
	\end{lstlisting}
	
	\subsubsection{strtok}
	\begin{lstlisting}
		// char *strtok(char *str, const char *delim);
		// Erster Aufruf: str; danach NULL. Modifiziert Original!
		char *tok = strtok(s, ",");
		while (tok != NULL) { tok = strtok(NULL, ","); }
	\end{lstlisting}
	
	\subsubsection{fflush / fclose / remove}
	\begin{lstlisting}
		fflush(f);           // Puffer sofort schreiben
		fclose(f);           // Puffer leeren + schliessen
		remove("data.txt");  // Datei loeschen
	\end{lstlisting}
	
	\subsubsection{fseek / ftell / rewind}
	\begin{lstlisting}
		// int fseek(FILE *stream, long offset, int whence);
		// whence: SEEK_SET, SEEK_CUR, SEEK_END
		fseek(f, 4, SEEK_SET);
		long size = ftell(f);
		rewind(f);  // == fseek(f, 0, SEEK_SET)
	\end{lstlisting}
	
	%---------------------------------------------------------------
	\section{Präprozessor \& Modulkonzept}
	
	\textbf{Build-Pipeline:} Präprozessor → Compiler (.o) → Linker → Loader
	
	\subsubsection{Direktiven}
	\begin{lstlisting}
		#include <stdio.h>       // Systemverzeichnis
		#include "datei.h"       // Projektverzeichnis
		
		#define ANZAHL 100
		#define QUADRAT(X) ((X)*(X))  // Klammern wichtig!
		// QUADRAT(3+5) ohne Klammern -> 3+5*3+5=23 statt 64
		
		#ifdef DEBUG
		printf("debug\n");
		#endif  // Aktivieren: gcc -DDEBUG ...
	\end{lstlisting}
	
	\subsubsection{Include Guard}
	\begin{lstlisting}
		#ifndef MEINE_DEFS_H
		#define MEINE_DEFS_H
		// Inhalt der Header-Datei
		#endif
	\end{lstlisting}
	
	\subsubsection{Modulkonzept}
	\begin{lstlisting}
		// doit.h
		#ifndef DOIT_H
		#define DOIT_H
		#define MAX 1000
		int doit(int b);
		#endif
		
		// doit.c
		#include "doit.h"
		int doit(int b) { return MAX + b; }
		
		// gcc -c main.c && gcc -c doit.c && gcc -o prog *.o
	\end{lstlisting}
	
	\subsubsection{Sichtbarkeit}
	\begin{lstlisting}
		int globalFunc(int a) { ... }        // ueberall sichtbar
		static int lokalFunc(int b) { ... }  // nur im Modul
	\end{lstlisting}
	
	\subsubsection{Prinizipien Modularisierung}
	\begin{tabular}{ll}
		\textbf{Geringe Vernetzung}  & wenige Schnittstellen \\
		\textbf{Schwache Koppplung} & schmale Schnittstellen \\
		\textbf{Hohe Kohäsion}  & inhaltliche Einheit - wohledfinierte Aufgabe \\
		\textbf{Explizite Schnittstelle} & Kommunikation mit Parameterübergrabe \\
		\textbf{Information Hiding} & nur global was muss; Verstecken Implementationsdetails; Trennung Interface \& Implementation \\
	\end{tabular}

	
	%---------------------------------------------------------------
	\section{Dynamische Speicherverwaltung}
	
	\begin{tabular}{ll}
		\textbf{Code}  & Programmcode \\
		\textbf{Daten} & globale/statische Variablen \\
		\textbf{Heap}  & dynamisch alloziert \\
		\textbf{Stack} & lokale Variablen (temporär) \\
	\end{tabular}
	
	\begin{tcolorbox}
		Keine Adresse lokaler Variablen zurückgeben -- Stack wird nach \texttt{return} abgebaut!
	\end{tcolorbox}
	
	\subsubsection{Funktionen (\texttt{<stdlib.h>})}
	\begin{lstlisting}
		void *malloc(size_t size);             // anfordern
		void  free(void *ptr);                 // freigeben
		void *calloc(size_t num, size_t size); // anfordern + 0
		void *realloc(void *ptr, size_t size); // Groesse anpassen
	\end{lstlisting}
	Immer \texttt{NULL}-Check nach \texttt{malloc}; \texttt{ptr = NULL} nach \texttt{free}.
	
	\subsubsection{Beispiele}
	\begin{lstlisting}
		// Array
		int  *arr = malloc(n * sizeof(int));   
		free(arr); 
		arr=NULL;
		// Struct
		Zeit *pZ  = malloc(sizeof(Zeit));      
		pZ->h = 4; 
		free(pZ);
		
		// String: erst grosser Buffer, dann passend
		char *buf = malloc(256);
		fgets(buf, 256, stdin);
		char *s = malloc(strlen(buf) + 1);
		strcpy(s, buf);  free(buf);  free(s);
	\end{lstlisting}
	
	\subsubsection{2D-Arrays}
	\begin{lstlisting}
		// Variante 1: 1D, Index manuell
		int *a = malloc(rows * cols * sizeof(int));
		a[i*cols + j] = ...;
		
		// Variante 2: Array von Zeigern (a[i][j])
		int **a = malloc(rows * sizeof(int *));
		for (int i = 0; i < rows; i++)
		a[i] = malloc(cols * sizeof(int));
		for (int i = 0; i < rows; i++) free(a[i]);
		free(a);
	\end{lstlisting}
	
	\subsubsection{Typische Fehler}
	\begin{lstlisting}
		double *p = malloc(n * sizeof(n));      // FALSCH: sizeof(n)
		double *p = malloc(n * sizeof(double)); // korrekt
		Datum  *d = malloc(10 * sizeof(Datum*));// FALSCH: sizeof(Datum*)
		Datum  *d = malloc(10 * sizeof(Datum)); // korrekt
		
		int *ptr = malloc(...);
		ptr = &array[0];  // Memory Leak: Heap verloren!
	\end{lstlisting}
	
	%---------------------------------------------------------------
	\section{Verkettete Liste}
	
	\begin{tcolorbox}
		Arrays: Lücken beim Entfernen $\Rightarrow$ teures Umkopieren. Listen: individuelle Allokation pro Element.
	\end{tcolorbox}
	
	\begin{lstlisting}
		typedef struct ListElem_s {
			char name[16]; char first[16]; int rank;
			struct ListElem_s *next;
		} ListElem;
	\end{lstlisting}
	
	\subsubsection{Element einfügen}
	\begin{lstlisting}
		// Mittendrin (nach item): Reihenfolge wichtig!
		newItem->next = item->next;
		item->next    = newItem;
		
		// Am Anfang:
		newItem->next = head;  head = newItem;
		
		// Am Ende:
		while (item->next != NULL) item = item->next;
		item->next = newItem;  newItem->next = NULL;
	\end{lstlisting}
	
	\subsubsection{Element entfernen}
	\begin{lstlisting}
		// Nach item:
		delPtr = item->next;  item->next = delPtr->next;
		free(delPtr);
		
		// Am Anfang:
		delPtr = head;  head = head->next;  free(delPtr);
		
		// Am Ende:
		while (delPtr->next != NULL)
		{ delPtr = delPtr->next; item = item->next; }
		item->next = NULL;  free(delPtr);
	\end{lstlisting}
	
	\subsubsection{Liste ausgeben}
	\begin{lstlisting}
		for (ListElem *it = head; it != NULL; it = it->next)
		printf("%s %s\n", it->first, it->name);
	\end{lstlisting}
	
	\subsubsection{Varianten}
	\textbf{Ankerelement:} Dummy-Head eliminiert Sonderfall am Anfang.\\
	\textbf{Tail-Zeiger:} append O(1) statt O(n).\\
	\textbf{Doppelt verkettet:} \texttt{prev}-Zeiger für Navigation in beide Richtungen.
	
	\subsubsection{Anwendungen \& Fehler}
	\textbf{Queue (FIFO):} enqueue / dequeue an verschiedenen Enden.\\
	\textbf{Stack (LIFO):} push / pop am selben Ende.\\
	\textbf{Binärbaum:} left/right-Zeiger, Suche O(log n).
	
	\begin{itemize}
		\item Speicher zu früh freigeben $\Rightarrow$ Pointer verloren
		\item Falsche Reihenfolge beim Umhängen $\Rightarrow$ Liste abgetrennt
		\item \texttt{free()} vergessen $\Rightarrow$ Memory Leak
	\end{itemize}
	
	%---------------------------------------------------------------
	\section{Strings}
	\begin{lstlisting}
		char  c      = 'A';     // Char-Literal: 1 Zeichen, Typ char
		char *s      = "Hallo"; // String-Literal: Zeiger, read-only
		char  arr[64];          // Array: fixer Speicher, beschreibbar
	\end{lstlisting}
	
	\subsubsection{Zuweisung}
	\begin{lstlisting}
		char arr[] = "Max";      // OK -- bei Deklaration
		arr[0] = '\0';           // OK -- einzelnes Zeichen
		strcpy(arr, "Max");      // OK -- ganzen String kopieren
		arr    = "Max";          // (*@\textcolor{red}{FEHLER}@*) -- Array ist kein l-value
		arr[0] = "\0";           // (*@\textcolor{red}{FEHLER}@*) -- String-Literal != char
	\end{lstlisting}
	
	\begin{center}
		\begin{tabular}{lll}
			\toprule
			& \texttt{'\textbackslash0'} & \texttt{"\textbackslash0"} \\
			\midrule
			Typ    & \texttt{char} (1 B)  & \texttt{char*} \\
			Wert   & Zeichen-Wert         & Adresse \\
			Grösse & 1 Byte               & 2 Bytes \\
			\bottomrule
		\end{tabular}
	\end{center}
	
	\begin{tabular}{ll}
		\toprule
		Situation & Typ \\
		\midrule
		Struct-Member              & \texttt{char name[N]} \\
		Dynamische Länge           & \texttt{char*} + malloc \\
		Funktionsparameter (lesen) & \texttt{const char*} \\
		\bottomrule
	\end{tabular}
	
	\begin{tcolorbox}
		\texttt{char name[N]} wann immer möglich -- kein malloc, kein free, keine Leaks.
	\end{tcolorbox}
	
	\begin{lstlisting}
		typedef struct { char first[64]; char name[64]; } Participant;
		memset(p, 0, sizeof(Participant));  // alles auf 0/NULL
	\end{lstlisting}
	
	%---------------------------------------------------------------
	\section{Arrays}
	\begin{lstlisting}
		int a[5];                         // uninitialisiert
		int a[5] = {1, 2, 3, 4, 5};      // vollstaendig
		int a[]  = {1, 2, 3, 4, 5};      // Groesse auto
		int a[5] = {1, 2};                // Rest = 0
		int a[5] = {0};                   // alles 0
		int *a   = (int[]){1,2,3};        // Compound Literal (C99)
		
		a[2] = 42;   int x = a[2];   *(a+2) = 42;  // aequivalent
	\end{lstlisting}
	
	\subsubsection{Arrays \& Pointer}
	\texttt{a == \&a[0]}; \texttt{a[i]} $\equiv$ \texttt{*(a+i)}
	\begin{lstlisting}
		int *p = a;  p[2];  *(p+2);  // alle gleich wie a[2]
	\end{lstlisting}
	
	\subsubsection{Mehrdimensional}
	\begin{lstlisting}
		int m[3][4] = {{1,2,3,4},{5,6,7,8},{9,10,11,12}};
		m[1][2] = 42;
	\end{lstlisting}
	
	\subsubsection{Als Funktionsparameter}
	\begin{lstlisting}
		void foo(int *a, int n);    // als Pointer
		void foo(int a[], int n);   // aequivalent
		void foo(int a[10], int n); // Groesse wird ignoriert
	\end{lstlisting}
	\begin{tcolorbox}
		Arrays werden \textbf{immer als Pointer} übergeben -- Länge separat mitgeben.
	\end{tcolorbox}
	
	%---------------------------------------------------------------
	\section{Structs}
	\begin{lstlisting}
		typedef struct { char name[16]; int rank; } Person;
		
		Person p;                                  // uninitialisiert
		Person p = {"Muster", 1};                 // positionsbasiert
		Person p = {.name="Muster", .rank=1};     // designated init
	\end{lstlisting}
	
	\subsubsection{Member-Zugriff}
	\begin{lstlisting}
		p.rank = 1;         // direkt: . Operator
		p.name[0];          // Array-Member
		
		Person *ptr = &p;
		ptr->rank = 1;      // ueber Pointer: -> Operator
		(*ptr).rank = 1;    // aequivalent, unueblich
	\end{lstlisting}
	
	\begin{tcolorbox}
		\texttt{p.member} wenn \texttt{p} ein Struct ist.\\
		\texttt{p->member} wenn \texttt{p} ein Pointer ist $\equiv$ \texttt{(*p).member}
	\end{tcolorbox}
	
	\subsubsection{Funktionsparameter \& malloc}
	\begin{lstlisting}
		void foo(Person p);   // Kopie -- lokale Aenderungen
		void foo(Person *p);  // Pointer -- wirkt auf Original
		
		Person *p = malloc(sizeof(Person));
		p->rank = 1;  free(p);
	\end{lstlisting}
	
\end{document}